\section{Sampling Design}
\label{sec:sampling}

This section describes the construction of an observationally informed prior
over the environmental parameters that enter \TIGRESS\ as initial conditions,
and the quasi-random design that converts this prior into a concrete set of
simulation inputs. The design fields, their identification with \PHANGS\
aperture measurements, and the validation fields reserved for post-simulation
comparison are defined in \autoref{sec:data}; the mapping to \TIGRESS\ initial
conditions is detailed in \autoref{sec:tigress_mapping}.

%-----------------------------------------------------------------
\subsection{Statistical Goal}
\label{sec:sampling:goal}
%-----------------------------------------------------------------

The objective of the sampling procedure is to construct an informed prior over
the environmental design vector
\begin{equation}
\thetaenv \;=\; \rbrackets{\Sgas,\ \Sstar,\ \Hstar,\ \Omrot,\ \qshear},
\label{eq:theta_env}
\end{equation}
conditioned on a target gas surface density band,
\begin{equation}
p\rbrackets{\thetaenv \,\big|\, \log\Sgas \in
\sbrackets{\log\Sgas^{\rm tgt} \pm \Delta}}.
\label{eq:prior_definition}
\end{equation}
This prior is \emph{not} an exact resampling of the \PHANGS\ aperture
distribution. Exact resampling would inherit the survey's selection effects
(e.g.\ inclination, distance, sensitivity, aperture footprint) and would tile
parameter space inefficiently for emulator and simulation-based inference
(\SBI) training. Instead, the goal is a compact set of simulation inputs that:
(i) preserves the dominant observed correlations among the five design fields;
(ii) covers the relevant parameter volume sufficiently densely for emulator
training and posterior coverage tests; and (iii) excludes parameter
combinations that are unphysical or implausible.

We deliberately exclude the molecular gas fraction $\fmol$ and the star
formation surface density $\Ssfr$ from the design vector. Neither is a
\TIGRESS\ input parameter: in the \TIGRESS\ framework both are emergent
outputs of the explicit thermochemistry and self-consistent star formation
prescription \citep{2024ApJ...972...67K}.\CGK{Add a second TIGRESS-NCR reference here if desired (e.g.\ Kim et al.\ 2023 ApJS thermochemistry paper); the key 2023ApJS..264...10K is not yet in references.bib and should be verified in ADS before insertion.} Requiring valid
$\fmol$ and $\Ssfr$ measurements in the reference set would furthermore bias
the prior toward high-completeness environments in which CO and SFR tracers
are simultaneously detected, suppressing the very low-$\fmol$ and
low-$\Ssfr$ regions that the simulation suite must cover. These two fields
are instead retained as \emph{validation quantities} to be compared against
the corresponding emergent distributions from the completed simulation
suite (\autoref{sec:results}).

%-----------------------------------------------------------------
\subsection{Design Fields and Validation Fields}
\label{sec:sampling:fields}
%-----------------------------------------------------------------

We partition the relevant \PHANGS\ aperture-level quantities into two
disjoint sets: design fields, which serve as \TIGRESS\ initial conditions
and over which the prior is constructed, and validation fields, which are
emergent in \TIGRESS\ and used only for post-simulation comparison.

\paragraph{Design fields.}
The five-dimensional design vector $\thetaenv$ (defined in
\autoref{eq:theta_env}) contains the gas surface density $\Sgas$, the
stellar surface density $\Sstar$, the stellar scale height $\Hstar$,
the orbital frequency $\Omrot$, and the shear parameter
$\qshear \equiv -d\ln\Omrot/d\ln R$.
These five quantities are sufficient to specify the local shearing-box
\TIGRESS\ initial conditions up to the dark matter midplane density
$\rhodm$, which is derived from the rotation curve as described in
\autoref{sec:tigress_mapping}.

\paragraph{Validation fields.}
The validation set comprises emergent quantities that are diagnostic of
the multiphase ISM and star formation, but are not used to specify the
simulations:
\begin{equation}
\thetaenv^{\rm val} = \rbrackets{\fmol,\ \Satom,\ \Smol,\ \Ssfr,\ \PDE,\ \seff}.
\end{equation}
Among these, $\PDE$ and $\seff$ are derived quantities in the \PRFM\
framework \citep{2022ApJ...936..137O} that are computed from the design
fields via the vertical equilibrium and feedback closure relations;
their \PHANGS\ values therefore depend on the same five design fields
plus the choice of \PRFM\ effective velocity dispersion model. We use
the design-field/validation-field split exclusively to define which
distributions are matched in the prior construction (design fields)
versus checked in post-simulation comparison (validation fields).

%-----------------------------------------------------------------
\subsection{Reference Pixel Selection}
\label{sec:sampling:reference}
%-----------------------------------------------------------------

The reference dataset $\Dref$ from which the prior is fit is the subset of
\PHANGS\ apertures \citep{2022AJ....164...43S,2023ApJ...945L..19S} that
satisfy a target $\Sgas$ band and basic physical validity. Specifically,
a \PHANGS\ aperture $i$ is included in $\Dref$ if
\begin{align}
\bigl|\log\Sgas^{(i)} - \log\Sgas^{\rm tgt}\bigr| &\leq \Delta,
\label{eq:sgas_band}\\
\Sgas^{(i)},\ \Sstar^{(i)},\ \Hstar^{(i)},\ \Omrot^{(i)} &> 0,
\label{eq:positivity}\\
0 &< \qshear^{(i)} \leq 1.5.
\label{eq:q_bound}
\end{align}
The band width $\Delta$ defines the local extent of $\Dref$ in $\log\Sgas$
and is chosen to balance two competing requirements: $\Delta$ must be wide
enough to provide a statistically meaningful sample for the kernel density
estimator, yet narrow enough that the conditional joint distribution of the
remaining design fields is well approximated as locally stationary in
$\log\Sgas$. A fiducial choice $\Delta = 0.2\dex$ is adopted in this work
and discussed further in \autoref{sec:results}.

The upper bound on the shear parameter, $\qshear \leq 1.5$, follows from
the definition $\qshear = -d\ln\Omrot/d\ln R$: $\qshear=0$ corresponds to
solid-body rotation, $\qshear=1$ to a flat rotation curve, and
$\qshear=1.5$ to Keplerian rotation around a point mass. Values
$\qshear > 1.5$ would correspond to super-Keplerian rotation curves that
cannot be produced by a positive enclosed mass distribution and are
therefore excluded as unphysical. We similarly exclude $\qshear \leq 0$,
which corresponds to rotation curves that increase faster than linearly
with radius and is not a regime in which the local shearing-box
approximation underlying \TIGRESS\ is expected to apply.

Crucially, validity of $\fmol$, $\Ssfr$, or any other tracer-specific
quantity is \emph{not} required for inclusion in $\Dref$. \CGK{This choice
maximizes the number of usable reference pixels — particularly in the
low-$\Sgas$, low-$\fmol$ outer-disk regime where CO is undetected and the
PHANGS megatable returns upper limits on $\Smol$ — but should be verified
against the megatable flag convention to ensure that apertures with all
five design fields validly measured are retained.}

%-----------------------------------------------------------------
\subsection{Kernel Density Estimator Prior}
\label{sec:sampling:kde}
%-----------------------------------------------------------------

We work throughout in log-transformed design space. Define the
five-dimensional log-vector
\begin{equation}
\bm{w} \;\equiv\;
\rbrackets{\log\Sgas,\ \log\Sstar,\ \log\Hstar,\ \log\Omrot,\ \log\qshear},
\label{eq:w_def}
\end{equation}
so that the reference set in log-space is $\cbrackets{\bm{w}_i}_{i=1}^{N}$
with $N = |\Dref|$. The log transformation is natural because each design
field spans more than an order of magnitude across \PHANGS\ and is
approximately log-normally distributed within fixed $\Sgas$ bands; it also
makes the KDE bandwidth selection scale-invariant under multiplicative
rescaling of any field.

We fit a Gaussian kernel density estimator to the reference log-vectors,
\begin{equation}
\hat{p}(\bm{w}) \;=\; \frac{1}{N}\,\sum_{i=1}^{N}
K_{\bm{H}}\rbrackets{\bm{w} - \bm{w}_i},
\label{eq:kde}
\end{equation}
where $K_{\bm{H}}$ is the multivariate Gaussian kernel
\begin{equation}
K_{\bm{H}}(\bm{u}) \;=\;
\rbrackets{2\pi}^{-5/2}\,\det(\bm{H})^{-1/2}\,
\exp\!\sbrackets{-\tfrac{1}{2}\,\bm{u}^{\top}\bm{H}^{-1}\bm{u}},
\label{eq:gaussian_kernel}
\end{equation}
and $\bm{H}$ is a symmetric positive-definite bandwidth matrix. We adopt
the multivariate Scott's rule for the full bandwidth matrix,
\begin{equation}
\bm{H} \;=\; \left(\frac{4}{d+2}\right)^{\!\!2/(d+4)}
\!\! N^{-2/(d+4)}\,\hat{\bm{\Sigma}},
\label{eq:scott_rule}
\end{equation}
where $d = 5$ is the number of design dimensions,
$N = |\Dref|$ is the reference sample size, and $\hat{\bm{\Sigma}}$ is
the empirical covariance of the $N$ reference log-vectors; for $d=5$
this gives $\bm{H} = (4/7)^{2/9} N^{-2/9} \hat{\bm{\Sigma}}$.
Multiplying the full empirical covariance by a scalar factor ensures
that $\bm{H}$ is positive definite and correctly encodes the observed
correlations among all five log-fields. In practice we will compare
against likelihood-based cross-validation to verify that this bandwidth
does not oversmooth the joint distribution.
Because all five design fields are strictly positive in $\Dref$ and
their upper bounds (most importantly $\qshear \leq 1.5$) are handled
by the post-mapping rejection step rather than by a boundary-corrected
KDE, the KDE in log-space requires no special boundary treatment.
This is in contrast to designs that include the molecular fraction
$\fmol \in [0,1]$, where a logit transform is needed to avoid
density leakage outside the physical support.

From $\hat{p}(\bm{w})$ we draw a large auxiliary KDE sample of size
$M \gg N$ (we use $M = 10^{5}$) and from this sample compute, for each
design coordinate $j = 1,\ldots,5$, the marginal cumulative distribution
function $F_j(w_j)$ and its inverse, the quantile function
\begin{equation}
Q_j(u) \;=\; F_j^{-1}(u), \qquad u \in (0,1).
\label{eq:quantile_function}
\end{equation}
The quantile functions are tabulated on a dense grid in $u$ and used as
look-up tables in the Sobol mapping (\autoref{sec:sampling:sobol}).
We emphasize that $\hat{p}(\bm{w})$ retains the full multivariate
correlation structure through the off-diagonal entries of $\bm{H}$,
but the marginal quantile functions $\cbrackets{Q_j}$ encode only the
one-dimensional marginals of $\hat{p}$. The dimension-wise
inverse-CDF mapping in \autoref{eq:sobol_map} therefore produces a
sample whose one-dimensional marginal distributions match those of
$\hat{p}$ \emph{exactly}, while the joint distribution corresponds to
an independence copula applied to these marginals. The joint
correlation structure of $\hat{p}$ is thus reproduced only
approximately, with fidelity depending on how well an independence
copula approximates the true KDE joint. In practice this approximation
is adequate when the pairwise correlations among the log design fields
are moderate, as is typical within a fixed $\Sgas$ band; we assess the
quality of the joint match quantitatively in \autoref{sec:results}.

%-----------------------------------------------------------------
\subsection{Sobol Sequence Design}
\label{sec:sampling:sobol}
%-----------------------------------------------------------------

Having reduced the prior to its marginal quantile functions on a unit
hypercube, we map a quasi-random low-discrepancy sequence through these
functions to produce the simulation design. We adopt a scrambled Sobol
sequence \citep{Sobol1967,Owen1998} rather than a pseudo-random or Latin
Hypercube Sampling (LHS) design for two reasons. First, for any finite
sample size $n$, Sobol sequences achieve a star-discrepancy of
$O(n^{-1}(\log n)^{d})$ in $d$ dimensions, asymptotically smaller than
the $O(n^{-1/2})$ scaling of pseudo-random samples and uniformly better
than LHS in projections onto coordinate subspaces. Second, and more
importantly for an incrementally grown simulation suite, the Sobol
construction has a nested progressive property: the first $2^{k}$ points
of the sequence form a balanced design at all $k$, so doubling the sample
size from $2^{k}$ to $2^{k+1}$ adds new design points without displacing
any of the original ones. LHS, by contrast, must in general be
re-stratified — and therefore re-drawn from scratch — whenever the
sample size changes.

\paragraph{Sampling algorithm.}
Given a target sample size $n$, the design is generated as follows.
\begin{enumerate}
\item Draw a scrambled Sobol sequence
$\cbrackets{\bm{u}_k}_{k=1}^{n} \subset [0,1]^{5}$ using
\texttt{scipy.stats.qmc.Sobol(d=5, scramble=True)} with a fixed random
seed for reproducibility.
\item Map each Sobol point through the marginal quantile functions of
\autoref{eq:quantile_function} to obtain log-space design points,
\begin{equation}
w_{k,j} \;=\; Q_j\rbrackets{u_{k,j}},
\qquad j = 1,\ldots,5.
\label{eq:sobol_map}
\end{equation}
\item Convert to physical-space design points,
\begin{equation}
\theta_{k,j} \;=\; 10^{\,w_{k,j}},
\label{eq:to_physical}
\end{equation}
yielding $\bm{\theta}_k = (\Sgas, \Sstar, \Hstar, \Omrot, \qshear)_k$.
\item For each design point with $\qshear_k > 1.5$, draw the next
un-used Sobol point from the sequence and repeat steps 2--3 until an
accepted point is obtained. Let $n_{\rm extra}$ denote the total
number of rejected-and-replaced draws. The accepted set of $n$ points
is formed from the first $n + n_{\rm extra}$ Sobol indices, with the
$n_{\rm extra}$ rejected indices recorded so that future extensions
begin from index $n + n_{\rm extra} + 1$. This bookkeeping preserves
the progressive nesting property (\autoref{sec:sampling:progressive}).
\end{enumerate}
The shear bound is the only post-mapping rejection criterion; positivity
is automatically guaranteed by the log-space construction. Because the
$\qshear$ marginal in $\Dref$ is already upper-truncated by
\autoref{eq:q_bound}, the rejection fraction $n_{\rm extra}/n$ is
expected to be small; we report the realized value in
\autoref{sec:results}.

Sample sizes that are powers of two, $n = 2^{m}$ with $m$ a non-negative
integer, are preferred because the Sobol sequence's discrepancy bounds
are sharp at these sizes and the construction is balanced in projections
to dyadic subintervals. Arbitrary $n$ is permitted at the cost of a mild
loss of projection uniformity.

%-----------------------------------------------------------------
\subsection{Progressive Extension of the Suite}
\label{sec:sampling:progressive}
%-----------------------------------------------------------------

A central operational advantage of the Sobol-based design is that the
suite is extensible at any later stage without re-drawing the existing
simulations. Let $S(n)$ denote the set of accepted design points produced by the
algorithm of \autoref{sec:sampling:sobol} at sample size $n$.
The underlying Sobol sequence has the nesting property: the first
$2^{k}$ sequence points are a subset of the first $2^{k+1}$ sequence
points. Provided the same random seed and the same physical rejection
criterion ($\qshear > 1.5$) are applied at every stage, the
accepted-point sets inherit this property,
\begin{equation}
S(2^{k}) \;\subset\; S(2^{k+1}) \;\subset\; S(2^{k+2}) \;\subset\; \cdots,
\label{eq:sobol_nesting}
\end{equation}
so that any simulation completed at stage $k$ is automatically a member
of the design at every subsequent stage. In practice we generate
$n + n_{\rm extra}$ Sobol points at each stage, accept the first $n$
points that pass the rejection criterion, and record the consumed Sobol
indices so that future extensions begin from the next un-used point in
the same sequence. The first generation of the
suite targets $n \in \{64, 128, 256\}$ ($m \in \{6, 7, 8\}$),
corresponding to a sequence of doublings; this progression is set by the
available computational allocation rather than any statistical
threshold, and the discrepancy of the design improves monotonically
along the sequence.

At each stage, the new $2^{k}$ points fill the largest remaining gaps in
the previous-stage coverage of the unit hypercube and, by the inverse-CDF
mapping of \autoref{eq:sobol_map}, the largest remaining gaps in
log-design space weighted by the KDE prior. This contrasts sharply with
LHS-based designs, which must be re-stratified to add new points and
therefore either displace the existing design or accept a strictly
sub-optimal stratification on the enlarged design.

This progressive structure is closely analogous to the staged extension
strategy adopted by the \CAMELS\ project
\citep{2021ApJ...915...71V}, in which cosmological hydrodynamic
simulation suites have been incrementally enlarged as new scientific
questions are identified. The Sobol nesting provides for the
\PHANGS-informed \TIGRESS\ suite the same operational property that
\CAMELS\ achieves through its explicit batch design: completed
simulations are never invalidated by the addition of new ones.

%-----------------------------------------------------------------
\subsection{Block-Structured Extension to Physics Parameters}
\label{sec:sampling:physics}
%-----------------------------------------------------------------

The first-generation suite holds all physics parameters fixed at fiducial
solar-neighborhood values (\autoref{sec:tigress_mapping}). Subsequent
generations will vary uncertain physics parameters such as the gas-phase
metallicity $\Zg$, the dust-to-gas ratio $\Zd$, the cosmic-ray
ionization rate $\xicr$, the assumed initial mass function, and
parameters of the supernova feedback prescription. We collect these
into a $d_{\rm phys}$-dimensional physics design vector $\thetaphys$.

We adopt a block-structured Sobol design over $(\thetaenv, \thetaphys)$,
defined as follows. Let
\begin{equation}
\cbrackets{\bm{u}^{\rm env}_k}_{k=1}^{n_{\rm env}}
\;\subset\; [0,1]^{5},
\qquad
\cbrackets{\bm{u}^{\rm phys}_\ell}_{\ell=1}^{n_{\rm phys}}
\;\subset\; [0,1]^{d_{\rm phys}}
\label{eq:two_sobol}
\end{equation}
be two independent scrambled Sobol sequences, one over the environmental
block and one over the physics block. The combined design is constructed by index-aligned pairing:
simulation $k$ uses environmental point $\bm{u}^{\rm env}_k$ and
physics point $\bm{u}^{\rm phys}_k$, so the total number of
simulations equals $\min(n_{\rm env}, n_{\rm phys})$. Each block is
mapped through its respective prior: $\bm{u}^{\rm env}_k$ through the
KDE-based quantile functions of \autoref{eq:sobol_map}, and
$\bm{u}^{\rm phys}_k$ through prior quantile functions appropriate to
each physics parameter (e.g.\ log-uniform in $\Zg$ over a chosen range).
When the two blocks differ in size, the smaller block is extended by
drawing additional Sobol points before pairing.

The block construction has two advantages. First, the two blocks may be
extended independently: doubling $n_{\rm env}$ at fixed $n_{\rm phys}$
yields a denser environmental sampling at the original set of physics
points, while doubling $n_{\rm phys}$ at fixed $n_{\rm env}$ adds physics
variations to the existing environmental design. Both extensions inherit
the nesting property of \autoref{eq:sobol_nesting} within their
respective blocks. Second, the two priors are kept manifestly
independent, which is desirable when the physics-parameter prior is
itself uncertain or under active revision: an update to $p(\thetaphys)$
does not require re-drawing the environmental design.

This block structure is analogous to the separation between
``cosmological'' and ``astrophysical'' parameter suites in the \CAMELS\
program \citep{2021ApJ...915...71V}, which treats the two parameter
groups as independent design axes that can be extended along either axis
without re-running the other.

\paragraph{Option A: pre-allocated joint design.}
If the full set of varied physics parameters and their priors are known
\emph{a priori}, an alternative is to draw a single
$(5+d_{\rm phys})$-dimensional Sobol sequence over the joint space and
map each block through its respective prior. This yields the optimal
joint discrepancy at all sample sizes and preserves the nesting property
of \autoref{eq:sobol_nesting} in the joint space, but at the cost of
fixing $d_{\rm phys}$ from the outset: adding a new physics parameter
later requires re-drawing the design. We therefore prefer the
block-structured construction for the staged suite considered here, and
reserve the pre-allocated joint design for a future generation in which
the physics-parameter list is finalized.

%-----------------------------------------------------------------
\subsection{Expanded Prior for Rare Environments}
\label{sec:sampling:expanded}
%-----------------------------------------------------------------

The KDE prior of \autoref{sec:sampling:kde} faithfully reproduces the
\PHANGS\ design-field distribution within the target $\Sgas$ band, but
by construction it under-samples regions of parameter space that are
under-represented in \PHANGS. Such regions include, for example, very
low orbital frequencies $\Omrot$ characteristic of low-surface-brightness
or outer-disk environments, large stellar scale heights $\Hstar$
characteristic of thickened or perturbed disks, and shear parameters
$\qshear$ near the rising-rotation-curve ($\qshear \to 0$) or Keplerian
($\qshear \to 1.5$) extremes. For emulator and \SBI\ training it is
desirable to extend coverage into these tails to guard against
extrapolation in posterior inference.

We implement such extension in two complementary ways. (i) Broadening
the KDE bandwidth, $\bm{H} \to s\,\bm{H}$ with $s > 1$, inflates the
prior in all directions while preserving the empirical mean and
correlation structure; this is sufficient for modest extrapolation.
(ii) Directional extension of the acceptance bounds in selected
log-space directions, implemented by augmenting the quantile functions
$Q_j(u)$ at the tails (i.e.\ extending $Q_j$ for $u$ near $0$ or $1$
beyond the support of $\Dref$), supports targeted extension into rare
regimes such as low-$\Omrot$ environments. Expanded-prior simulations
are flagged at draw time and analyzed separately from the
\PHANGS-prior simulations in the fairness diagnostics
(\autoref{sec:sampling:fairness}) and downstream emulator validation.
\CGK{The precise broadening factor $s$ and the directional extensions
applied in the first-generation suite will be reported in
\autoref{sec:results}; the choice of these is closely tied to which
PHANGS-poor environments the inference program ultimately targets.}

%-----------------------------------------------------------------
\subsection{Sampling Fairness Diagnostic}
\label{sec:sampling:fairness}
%-----------------------------------------------------------------

We define a quantitative diagnostic of how faithfully the simulation
design reproduces the \PHANGS\ design-field distribution within the
target $\Sgas$ band. For each design field $a \in \thetaenv$, let
$\hat{q}^{\rm sample}_{a}(\tau)$ and $\hat{q}^{\rm ref}_{a}(\tau)$
denote the empirical $\tau$-th quantile of $\log_{10} a$ in the
simulation design and in $\Dref$, respectively (we use $\tau$ for
quantile probability levels to distinguish from the KDE marginal
quantile functions $Q_j$ of \autoref{eq:quantile_function}).
The per-field quantile mismatch is
\begin{equation}
\varepsilon_a \;=\;
\max_{\tau \in \cbrackets{0.1,\,0.2,\,\ldots,\,0.9}}
\Bigl|\,\hat{q}^{\rm sample}_{a}(\tau)
- \hat{q}^{\rm ref}_{a}(\tau)\,\Bigr|,
\label{eq:eps_a}
\end{equation}
in units of dex, and the overall design fairness is the worst-field
mismatch,
\begin{equation}
\varepsilon_{\rm env} \;=\;
\max_{a \in \thetaenv}\,\varepsilon_a.
\label{eq:eps_env}
\end{equation}
A design is considered \emph{fair} with respect to $\Dref$ if
$\varepsilon_{\rm env} \leq \varepsilon_{\rm tol}$, where we adopt
$\varepsilon_{\rm tol} = 0.10$--$0.20\dex$ as a fiducial tolerance.
This range is motivated by the typical aperture-level measurement
uncertainty in the underlying \PHANGS\ design fields, which is
itself of order $\sim 0.1\dex$ for the surface densities and somewhat
larger for $\Hstar$; a design that matches the reference marginal
quantiles to better than the measurement uncertainty is not
meaningfully distinguishable from the reference within $\Dref$.
\CGK{The 0.10--0.20\,dex tolerance is a working choice; a
quantitative justification from the PHANGS measurement-uncertainty
budget should be added when those numbers are tabulated.}

The diagnostic in \autoref{eq:eps_env} applies exclusively to the
five \emph{design} fields. The validation fields $\fmol$, $\Satom$,
$\Smol$, $\Ssfr$, $\PDE$, and $\seff$ are not used to define the
prior, and they cannot be matched at design time because they are
not \TIGRESS\ inputs. Instead, the simulation outputs for the
validation fields will be compared against the corresponding
\PHANGS\ distributions in \autoref{sec:results} using analogous
quantile-mismatch metrics, providing a post-hoc verification that
the simulated ISM populates the same emergent regime as the
observed one. A successful program is one in which
$\varepsilon_{\rm env} \leq \varepsilon_{\rm tol}$ by design and the
corresponding validation-field mismatches are also small without
having been tuned.
